\subsection*{7.4 Push-Forward Measure and Change-of-Variable Formula}

When modeling a stochastic experiment, we often assume that a sample is drawn from a probability space according to a probability measure. However, the probability space may be complex, and evaluating a Lebesgue integral on it may not always be straightforward. In some cases, we cannot directly observe the sample drawn from the sample space, but instead we measure the experiment using random variables.

To calculate the expectation of these random variables, it is often more convenient to transform the original probability space to a more concrete probability space defined on the real number line. If the Stieltjes measure function is differentiable, we can simplify a Lebesgue integral on the real line by transforming it to a Riemann integral. This simplification is useful, as Riemann integrals are often easier to compute than Lebesgue integrals. The reduction process is shown as follows:
\[
\text{Lebesgue integral on abstract prob. space}
\to
\text{Lebesgue--Stieltjes integral on } \mathbb{R}
\to
\text{Riemann--Stieltjes integral on } \mathbb{R}
\to
\text{Riemann integral}.
\]

In this section we consider an abstract measure space $(\Omega,\mathcal{F},\mu)$ and a real-valued measurable function $X$ that maps out of it.
\begin{equation}
(\Omega,\mathcal{F},\mu) \xrightarrow{X} (\mathbb{R},\mathcal{B}(\mathbb{R})).
\tag{7.6}
\end{equation}

\begin{thmbox}{7.10}
Given a measurable mapping $X$ as in (7.6), the set function
\[
X_{\#}\mu(B)\triangleq \mu(X^{-1}(B))=\mu(\{\omega : X(\omega)\in B\})
\]
is a Borel measure defined for all $B\in \mathcal{B}(\mathbb{R})$.
\end{thmbox}

\textit{Proof} Since $X$ is $(\mathcal{F},\mathcal{B}(\mathbb{R}))$-measurable, the pre-image $X^{-1}(B)$ is $\mathcal{F}$-measurable. Hence $X_{\#}\mu$ is well-defined. We check the axioms of measure below. It is obvious that $X_{\#}\mu(\emptyset)=0$. For any sequence of mutually disjoint sets $B_1,B_2,\dots$ in $\mathcal{B}(\mathbb{R})$,
\[
X_{\#}\mu\Bigl(\biguplus_{i=1}^{\infty} B_i\Bigr)
=
\mu\Bigl(X^{-1}\bigl(\biguplus_{i=1}^{\infty} B_i\bigr)\Bigr)
=
\mu\Bigl(\biguplus_{i=1}^{\infty} X^{-1}(B_i)\Bigr)
=
\sum_{i=1}^{\infty} X_{\#}\mu(B_i).
\]

The second equality follows from (4.3) and the third equality from the assumption that $\mu$ is a measure function. \hfill $\square$

\begin{defbox}{7.2}
The measure $X_{\#}\mu$ defined in Theorem 7.10 is called the \textit{push-forward measure} defined by $X$. It is also called the \textit{image} of $\mu$ or the \textit{measure induced by $X$}. When $\mu$ is a probability measure, the induced measure $X_{\#}\mu$ is also a probability measure and is called the \textit{distribution} of $X$.

Other notation for the push-forward measure include $X_{*}\mu$, $\mu^X$, $\mu_X$, and $\mu\circ X^{-1}$.
\end{defbox}

\textbf{Example 7.4.1 (Rayleigh Distribution)} \\
We generate a random point on $\mathbb{R}^2$, so that the $x$- and $y$-coordinates are independent and identically distributed Gaussian random variables. Suppose the mean is 0 and the variance is $\sigma^2$. We want to find the distribution of the distance of the random point to the origin.

We can formulate this problem by regarding $(\mathbb{R}^2,\mathcal{B}(\mathbb{R}),P)$ as the probability space and specify the probability measure $P$ on the semi-infinite rectangles
\[
P((-\infty,x]\times (-\infty,y])
=
\int_{-\infty}^{x} \frac{e^{-x^2/(2\sigma^2)}}{\sqrt{2\pi\sigma^2}}\, dx
\int_{-\infty}^{y} \frac{e^{-y^2/(2\sigma^2)}}{\sqrt{2\pi\sigma^2}}\, dy.
\]

The measure $P$ will then be uniquely defined by the measure extension theorem, as the Borel algebra on $\mathbb{R}^2$ is generated by the semi-infinite rectangles.

Let $Z$ denote the distance between the random point and the origin. The distance $Z$ is less than $z$ if and only if $X^2+Y^2\le z^2$. Hence, we can derive the cdf of $Z$ by
\[
P(Z\le z)
=
P(X^2+Y^2\le z^2)
=
\frac{1}{2\pi\sigma^2}
\iint_{x^2+y^2\le z^2}
e^{-(x^2+y^2)/(2\sigma^2)}\, dx\, dy.
\]

After some calculus, we can simplify it to $P(Z\le z)=1-e^{-z^2/(2\sigma^2)}$. The pushforward measure $Z_{\#}P$ is the probability measure on $\mathbb{R}$ induced by the Stieltjes measure function $F(z)=1-e^{-z^2/(2\sigma^2)}$, which is the cdf of a Rayleigh distribution.

The change-of-variable formula gives the relation between a Lebesgue integral on an abstract space and a Lebesgue integral on the real number line.

\begin{thmbox}{7.11 (Change-of-Variable Formula)}
Let $(\Omega,\mathcal{F},\mu)$ be a measure space, and consider two measurable maps $X$ and $h$
\[
(\Omega,\mathcal{F},\mu)\xrightarrow{X} (\mathbb{R},\mathcal{B}(\mathbb{R})) \xrightarrow{h} (\mathbb{R},\mathcal{B}(\mathbb{R})).
\]

Let $X_{\#}\mu$ denote the measure on $\mathbb{R}$ induced by $X$. Then,
\begin{enumerate}[label=(\alph*)]
    \item $h(X)\in L^1(\mu)$ iff $h\in L^1(X_{\#}\mu)$.
    \item If $h$ is nonnegative or if the two equivalent conditions in (a) hold, then
    \begin{equation}
    \int_{\Omega} h(X(\omega))\, d\mu(\omega)
    =
    \int_{\mathbb{R}} h(x)\, dX_{\#}\mu(x).
    \tag{7.7}
    \end{equation}
\end{enumerate}
\end{thmbox}

The function $h$ in the theorem may be continuous or piece-wise continuous. The important requirement is that the composite function $h(X)$ is measurable, so that the composite function $h\circ X$ is also measurable.

\textit{Proof} Suppose $h$ is an indicator function $1_B$ for some Borel set $B\in \mathcal{B}(\mathbb{R})$. The left-hand side of (7.7) is
\[
\int_{\Omega} h(X(\omega))\, d\mu(\omega)
=
\int_{\Omega} (1_B\circ X)(\omega)\, d\mu(\omega)
=
\int_{\Omega} 1_{X^{-1}(B)}\, d\mu
=
\mu(X^{-1}(B)).
\]

Meanwhile, the right-hand side is
\[
\int_{\mathbb{R}} 1_B\, dX_{\#}\mu = X_{\#}\mu(B)=\mu(X^{-1}(B)).
\]

The first equality above follows from the definition of Lebesgue integral for simple function and the second from the definition of $X_{\#}\mu$. Therefore (7.7) holds for indicator functions $h$.

Suppose $h$ is a nonnegative and measurable function. Let $h_n$ be simple and nonnegative functions for $n=1,2,3,\dots$ such that $h_n \nearrow h$. We have $h_n(X)\nearrow h(X)$. Apply the monotone convergence theorem two times in the following derivation:
\[
\int_{\Omega} h(X)\, d\mu
=
\int_{\Omega} \lim_{n\to\infty} h_n(X(\omega))\, d\mu(\omega)
\]
\[
\stackrel{\text{MCT}}{=}
\lim_{n\to\infty} \int_{\Omega} h_n(X(\omega))\, d\mu(\omega)
\]
\[
=
\lim_{n\to\infty} \int_{\mathbb{R}} h_n(x)\, dX_{\#}\mu(x)
\]
\[
\stackrel{\text{MCT}}{=}
\int_{\mathbb{R}} \lim_{n\to\infty} h_n(x)\, dX_{\#}\mu(x)
\]
\[
=
\int_{\mathbb{R}} h(x)\, dX_{\#}\mu(x).
\]

Therefore (7.7) holds for nonnegative and measurable functions.

Next we assume that $h$ is a real-valued and measurable function. Applying the argument in the previous paragraph to $|h|$, we obtain part (a) immediately.

Write $h=h^+-h^-$, where $h^+$ and $h^-$ are the positive and negative parts of $h$. Suppose $h(X)$ is integrable, i.e., suppose that $\int h^+(X)<\infty$ and $\int h^-(X)<\infty$. We note that $h^+(X)$ and $h^-(X)$ are both nonnegative and measurable functions in $L^1(\mu)$. The integral of $h(X)$ is
\[
\int_{\Omega} h(X)\, d\mu
\triangleq
\int_{\Omega} h^+(X(\omega))\, d\mu(\omega)
-
\int_{\Omega} h^-(X(\omega))\, d\mu(\omega)
\]
\[
=
\int_{\mathbb{R}} h^+(x)\, dX_{\#}\mu(x)
-
\int_{\mathbb{R}} h^-(x)\, dX_{\#}\mu(x)
\]
\[
=
\int_{\mathbb{R}} h\, dX_{\#}\mu.
\]

This finishes the proof of the change-of-variable formula. \hfill $\square$

When applied to finite probability spaces, the change-of-variable formula says that we can compute a function of a random variable $X$ by
\[
E[g(X)] = \sum_{i=0}^{\infty} g(i)p_i,
\]
where $p_i$ is the probability $P(X=i)$, for $i=0,1,2,\dots$, and $g$ is a real-valued function.

In the next theorem, the formula given in Eq.~(7.8) has earned the nickname ``LOTUS'' (Law of the Unconscious Statistician). Additionally, the formula in Eq.~(7.9) serves as the definition of the expectation of a continuous random variable in a first course in probability. We thus see that the measure-theoretic definition of expectation unifies the notion of mathematical expectation for discrete and continuous random variables.

\begin{thmbox}{7.12}
Let $X$ be a random variable defined on $(\Omega,\mathcal{F},P)$. Suppose the induced measure $X_{\#}P$ on $\mathbb{R}$ has Stieltjes measure function $F_X(x)$, and suppose $F_X(x)$ is differentiable with derivative $f_X(x)$. Then, for any piece-wise continuous function $g(x)$ such that $g(X)$ is $P$-integrable, $E[g(X)]$ is given by
\begin{equation}
E[g(X)] = \int_{-\infty}^{\infty} g(x)f_X(x)\, dx.
\tag{7.8}
\end{equation}

In particular, when the function $g(x)=x$ is the identity function, we recover the formula for computing the mean of a continuous random variable
\begin{equation}
E[X] = \int_{-\infty}^{\infty} x f_X(x)\, dx.
\tag{7.9}
\end{equation}
\end{thmbox}

\textit{Proof} By the change-of-variable formula in (7.7), the expectation of $g(X)$ is equal to the Lebesgue--Stieltjes integral
\[
\int_{\Omega} g(X)\, dP = \int_{\mathbb{R}} g(x)\, dX_{\#}P(x).
\]

We can transform it to a Riemann--Stieltjes integral
\[
\int_{-\infty}^{\infty} g(x)\, dF_X(x)
\]
by applying Theorem 7.9. Since $F_X(x)$ is differentiable, this can be further simplified to a Riemann integral $\int_{-\infty}^{\infty} g(x)f_X(x)\, dx$. \hfill $\square$

One immediate application of LOTUS is the computation of the mean and variance of a probability density distribution.

\textbf{Example 7.4.1 (Variance of a Random Variable)} \\
Let $X$ be a discrete random variable with pmf $P(X=i)=p_i$ for $i\ge 0$ and mean $m$. To evaluate the variance of $X$, we need not obtain the distribution of random variable $Y=(X-m)^2$, but rather compute it directly using the formula
\[
E[(X-m)^2] = \sum_i (i-m)^2 p_i.
\]

\textbf{Example 7.4.2 (Computing the Mean of a Chi Distribution)} \\
The chi distribution with $k$ degree of freedom is the square root of the chi-square distribution with the same degree of freedom. Although the pdf of the chi distribution is available in closed form, in the computation of the mean, however, we can simply use the pdf of chi-square distribution
\[
f_{\chi^2}(x)
=
\begin{cases}
\dfrac{1}{2^{k/2}\Gamma(k/2)}x^{k/2-1}e^{-x/2} & \text{if } x\ge 0,\\
0 & \text{if } x<0,
\end{cases}
\]
which is a special case of Gamma distribution. By Theorem 7.12, we can write
\[
\int_{0}^{\infty} \sqrt{x}\, f_{\chi^2}(x)\, dx
=
\int_{0}^{\infty}
\sqrt{x}\,
\frac{1}{2^{k/2}\Gamma(k/2)}x^{k/2-1}e^{-x/2}\, dx
\]
\[
=
\frac{\sqrt{2}\Gamma((k+1)/2)}{\Gamma(k/2)}
\int_{0}^{\infty}
\frac{1}{2^{(k+1)/2}\Gamma((k+1)/2)}
x^{\frac{k+1}{2}-1}e^{-x/2}\, dx.
\]

The integral is equal to $1$ because the summand is the pdf of a Gamma distribution. The expectation of the chi distribution with $k$ degree of freedom equals $\dfrac{\sqrt{2}\Gamma((k+1)/2)}{\Gamma(k/2)}$.
